\chapter{Grounding and Reliability}
\label{ch:grounding}

\epigraph{``An ungrounded agent is a hallucination engine with side effects.''}{}

\learninggoal{
	\item Define grounding and distinguish its dimensions (factual, contextual, procedural).
	\item Implement self-verification and consistency checks.
	\item Design human-in-the-loop feedback mechanisms.
	\item Understand the limits of grounding---what cannot be verified.
}

\section{The Grounding Problem}

Agents operate in the real world. When they act on incorrect information, the consequences are real. \textbf{Grounding} is the process of ensuring that an agent's beliefs, decisions, and actions are based on verifiable facts rather than model hallucination.

\begin{definition}[Grounding]
	An agent action is \textbf{grounded} if its justification is traceable to:
	\begin{enumerate}
		\item A trusted external source (database, API, document), or
		\item A verifiable chain of reasoning, or
		\item Explicit user confirmation.
	\end{enumerate}
	Actions that cannot satisfy any of these criteria are \textbf{ungrounded}.
\end{definition}

\subsection{Types of Grounding Failure}

\begin{longtable}{p{3.5cm}p{4.5cm}p{4.5cm}p{2.5cm}}
	\toprule
	\textbf{Failure type} & \textbf{Description}                & \textbf{Example}                  & \textbf{Detection} \\
	\midrule
	Factual               & Agent asserts false information     & ``The Eiffel Tower is in London'' & Cross-reference    \\
	Contextual            & Agent contradicts provided context  & Ignores retrieved documents       & Consistency check  \\
	Procedural            & Agent executes wrong procedure      & Calls wrong API with wrong args   & Schema validation  \\
	Planning              & Agent's plan cannot achieve goal    & Missing critical step             & Plan verification  \\
	Calibration           & Agent overconfident in wrong answer & High confidence, wrong result     & Confidence scoring \\
	\bottomrule
\end{longtable}

\section{Self-Verification}

The agent verifies its own outputs before acting:

\begin{lstlisting}[caption=Self-verification loop]
  def grounded_action(action, context):
      """Execute action only if verification passes."""
      # Step 1: Verify factual claims
      claims = extract_claims(action)
      for claim in claims:
          evidence = retrieve_evidence(claim)
          if not supports(evidence, claim):
              return reject(f"Unsupported claim: {claim}")

      # Step 2: Verify procedure
      if action.type == "tool_call":
          if not validate_schema(action, tool_registry):
              return reject("Invalid tool arguments")
          if not verify_prerequisites(action, context):
              return reject("Prerequisites not met")

      # Step 3: Verify plan coherence
      if action.type == "plan":
          verification = verify_plan(action, context.goal)
          if not verification.passes:
              return reject(verification.reason)

      return execute(action)
\end{lstlisting}

\subsection{Self-Consistency}

Sample multiple reasoning paths and check agreement:

\begin{equation}
	\text{consistency}(x) = \frac{2}{k(k-1)} \sum_{i<j} \mathbbm{1}\{f(x; \theta_i) = f(x; \theta_j)\}
\end{equation}

where $f$ is the agent's output and $\theta_i$ are independent samples (using temperature $\tau > 0$). Low consistency ($< 0.6$) indicates uncertainty.

\begin{principle}[Consistency Threshold]
	An agent should not execute an action with semantic consistency below 0.6 without human confirmation. The threshold is task-dependent: lower for reversible actions, higher for irreversible ones.
\end{principle}

\section{Human-in-the-Loop}

Some decisions are too important to trust to the agent alone. Human-in-the-loop (HITL) mechanisms provide a safety net.

\subsection{Approval Gates}

Sensitive actions require human approval before execution:

\begin{lstlisting}[caption=Approval gate flow]
  def execute_with_gate(action, user_channel):
      if action.requires_approval:
          prompt = f"""Approve this action?
          Action: {action.description}
          Arguments: {action.args}
          Consequences: {action.consequences}
          Reply 'approve' or 'reject'."""

          response = send_to_user(user_channel, prompt)
          if response != "approve":
              return reject("User rejected action")
      return execute(action)
\end{lstlisting}

Which actions should require approval?

\begin{itemize}
	\item \textbf{Always require approval:} irreversible actions (sending email, deleting files, making payments).
	\item \textbf{Conditional approval:} actions exceeding a cost/user/scope threshold.
	\item \textbf{Review after:} read-only actions that can be reverted (search, data retrieval).
\end{itemize}

\subsection{Escalation}

When the agent cannot resolve an ambiguity, it escalates:

\begin{lstlisting}[caption=Escalation logic]
  def should_escalate(state):
      reasons = []
      if len(state.uncertain_claims) > 3:
          reasons.append("Too many uncertain claims")
      if state.same_error_count > 2:
          reasons.append("Repeated failure on same action")
      if state.confidence < 0.3:
          reasons.append("Overall confidence too low")
      if time_since_last_confirmation > 10_minutes:
          reasons.append("Long operation without user check")
      return reasons

  def escalate(state, user_channel):
      summary = format_escalation_summary(state)
      send_to_user(user_channel, summary)
      response = wait_for_user_guidance()
      return response
\end{lstlisting}

\section{Tool Output Verification}

Tool outputs can be wrong, corrupted, or malicious. The agent must verify them:

\begin{lstlisting}[caption=Tool output verification]
  def verify_tool_output(output, expected_schema):
      """Check tool output for integrity and correctness."""
      checks = []

      # Schema validation
      try:
          parsed = json.loads(output)
          validate_schema(parsed, expected_schema)
          checks.append(("schema", True))
      except:
          checks.append(("schema", False))

      # Range checks (if applicable)
      if expected_schema.get("type") == "number":
          min_val = expected_schema.get("minimum", -inf)
          max_val = expected_schema.get("maximum", inf)
          in_range = min_val <= parsed <= max_val
          checks.append(("range", in_range))

      # Consistency with prior knowledge
      if known_facts := retrieve_related_facts(output):
          consistent = check_consistency(output, known_facts)
          checks.append(("consistency", consistent))

      return all(passed for _, passed in checks)
\end{lstlisting}

\section{Calibration}

An agent is \textbf{calibrated} if its confidence matches its reliability:

\begin{equation}
	\mathbb{P}(\text{action succeeds} \mid \text{confidence} = c) = c \quad \forall c \in [0, 1]
\end{equation}

Measuring calibration requires logging confidence and outcomes over many episodes:

\begin{equation}
	\text{ECE} = \sum_{m=1}^{M} \frac{|B_m|}{N} \left|\text{acc}(B_m) - \overline{\text{conf}}(B_m)\right|
\end{equation}

where $B_m$ are confidence bins (e.g., 0.0--0.1, 0.1--0.2, \dots).

If the agent is miscalibrated (common: overconfident in wrong answers), the confidence scores can be post-hoc calibrated using temperature scaling on the \llm's token probabilities.

\section{The Limits of Grounding}

Not everything can be grounded:

\begin{itemize}
	\item \textbf{Subjective judgments:} ``Is this design beautiful?'' cannot be verified against an external source.
	\item \textbf{Creative tasks:} ``Write a poem'' has no correctness criterion.
	\item \textbf{Future predictions:} ``Will it rain tomorrow?'' cannot be verified until tomorrow.
	\item \textbf{Open-ended exploration:} ``What are possible solutions?'' the agent must explore without grounding.
\end{itemize}

\begin{principle}[Grounding-Aware Design]
	An agent should know what it cannot ground and adjust its behaviour accordingly: express uncertainty, ask for confirmation, or flag actions as unverifiable.
\end{principle}

\section{Exercises}

\begin{exercise}
	Design a verification pipeline for a code-generation agent. For each generated function, what must be verified before execution? How do you handle dependencies (imports, libraries)?
\end{exercise}

\begin{exercise}
	Implement a self-consistency check for a factual QA agent. Use $k=5$ samples, temperature $\tau=0.7$. Measure consistency against accuracy on a test set of 100 questions.
\end{exercise}

\begin{exercise}
	Design an approval-gate policy for a personal finance agent (can read balances, schedule payments, and generate reports). Which actions are always gated, which are conditional, and which are unconstrained?
\end{exercise}

\begin{exercise}
	An agent's calibration curve shows $\text{acc}(0.9) = 0.6$---it is overconfident. Propose a post-hoc calibration method. How would you validate it?
\end{exercise}
